\section{Introduction}

\IEEEPARstart{R}{obotic} manipulators have become indispensable in modern manufacturing, industrial automation, medical robotics, and human–robot interaction due to their high precision, flexibility, and autonomous operation capabilities. In many robotic applications, accurate trajectory tracking is essential for ensuring motion precision, operational safety, and task efficiency. However, achieving high-performance tracking remains challenging because practical robotic systems are inevitably affected by modeling uncertainties, external disturbances, and unmodeled nonlinear dynamics. These uncertainties may significantly degrade tracking accuracy and even compromise the stability of conventional model-based control methods.

To improve tracking performance under uncertainties, a variety of nonlinear control approaches have been extensively investigated, including computed torque control, robust control, adaptive control, sliding mode control, and intelligent control methods. Among them, computed torque control has been widely adopted because of its systematic design methodology and excellent tracking performance when accurate dynamic models are available. Nevertheless, its effectiveness strongly depends on precise knowledge of the manipulator dynamics, which is often difficult to obtain in practice due to uncertain physical parameters, friction, payload variations, and external disturbances. Adaptive control alleviates part of this limitation by estimating unknown parameters online, while robust control improves disturbance rejection at the expense of increased control effort or conservative gain selection.

In recent years, neural-network-based adaptive control has attracted considerable attention because of its capability to approximate unknown nonlinear functions without requiring explicit mathematical models. In particular, radial basis function (RBF) neural networks have been widely employed owing to their universal approximation capability, simple network structure, and efficient online learning properties. By integrating RBF neural networks with adaptive control, unknown dynamics can be estimated online and compensated during controller execution, thereby improving tracking accuracy under uncertain environments. Nevertheless, maintaining parameter boundedness and ensuring stable online adaptation remain important issues in adaptive neural control, especially in the presence of approximation errors and persistent disturbances.

Motivated by these observations, this paper investigates the trajectory tracking problem for uncertain robotic manipulators using an adaptive neural control framework based on RBF networks. Unknown nonlinear dynamics are represented as lumped uncertainties and approximated online through RBF neural networks within a computed torque control architecture. Furthermore, a $\sigma$-modified adaptive law is incorporated to suppress parameter drift and enhance robustness during online learning. Stability of the overall control system is established through Lyapunov analysis, and numerical simulations are conducted to evaluate the effectiveness of the proposed approach.

The main contributions of this paper are summarized as follows.

An adaptive neural control framework combining computed torque control and RBF neural network approximation is developed for trajectory tracking of uncertain robotic manipulators.
A $\sigma$-modified adaptive law is incorporated into the online learning process to improve parameter robustness and suppress weight drift.
Lyapunov-based stability analysis and numerical simulations demonstrate the effectiveness of the proposed control strategy under model uncertainties and external disturbances.

The remainder of this paper is organized as follows. Section II introduces the preliminary knowledge of robotic manipulator dynamics and RBF neural networks. Section III formulates the trajectory tracking problem. Section IV presents the adaptive controller design together with the corresponding stability analysis. Section V provides numerical simulation results. Finally, Section VI concludes the paper.